\section{RELATED WORK}

In this section, we briefly review sparse motion tracking pipelines, standardized runtime I/O, and avatar framework for humanoid representation. 

\subsection*{XR Systems with Sparse Sensing}
Sparse-sensing estimation attempts to reconstruct full-body pose from a limited set of consumer inputs. Controller- and HMD-based regressors, such as AvatarPoser~\cite{jiang2022avatarposer}, show that plausible full-body motion can be inferred from only a headset and controllers, while robustness can be enhanced through auxiliary acceleration cues as in DivaTrack~\cite{yang2024divatrack}. 

Building on these foundations, recent work has advanced sparse-sensing through compact machine learning models with categorical matching~\cite{starke2024categorical}, and hardware-specific, scalable approaches leveraging commodity Inertial Measurement Unit (IMU) sensors.
For example, HMDPoser~\cite{dai2024hmd} addresses scalable HMD scenarios by optionally incorporating multiple physical trackers, all integrated within a single inference model.
This model is further enhanced by a temporal-spatial transformer module combined with SMPL estimation.
% Meanwhile, IMU-driven methods demonstrate practical deployment via widely available devices such as smartphones, smartwatches, and earbuds~\cite{xu2024mobileposer}. 

Despite these advances, practical challenges remain. 
Users struggle to select method that best match their requirements within a fragmented XR ecosystem, where platforms, hardware vendors, and even humanoid avatar formats differ.
As a result, choosing an appropriate method can be as difficult as achieving accurate pose reconstruction itself.

Moreover, perceptual outcomes depend not only on reconstruction accuracy but also on how motion is ultimately represented through avatar rigs and animation design choices~\cite{yun2023animationfidelity}. 
Enabling modular configurations that allow users to tune these factors is therefore essential for enhancing presence in XR environments.
Taken together, these findings suggest that sensing, completion, and presentation layers must be treated as modular and co-tunable, ensuring that technical advances align with user-level development.

\subsection*{Standardized Runtime I/O and Recent Cross-Platform Developments}
Despite growing adoption of cross-platform standards, runtime I/O in XR systems remains fragmented due to vendor-specific coordinate conventions, joint hierarchies, and update semantics. 
Vendors typically expose proprietary data layouts and frame orientations, requiring developers to reconcile heterogeneous inputs even within a single engine. 

OpenXR~\cite{khronos_openxr_1_0} and its extensions such as XR\_EXT\_hand\_tracking~\cite{openxr_ext_hand} have made significant strides in standardizing data acquisition interfaces. 
Even the Android XR platform and major vendors, including Meta, XREAL, and Vive, have committed to OpenXR compliance across their product lines~\cite{android_xr_sdk_2024,meta_and_openxr_2025,xreal_sdk_3_0_2024,vive_openxr_2024}, yet they continue to provide platform-specific data formats and legacy version support.
Furthermore, because OpenXR explicitly avoids prescribing end-to-end motion pipelines or pose semantics, integration of motion data into avatar representation remains the responsibility of the individual developer or user.
As a result, developers must still manually unify pose conventions, resolve transform mismatches, and reconfigure downstream processing logic for each supported SDK.

\subsection*{Avatar Frameworks for Humanoid Representation and Animation}
Based on the SMPL family~\cite{loper2015smpl}, which provides a learned, parameterized surface model of the human body, Meshcapade offers a comprehensive asset pipeline, API suite, and tooling infrastructure that serve as the representation backbone for many modern motion analysis and synthesis systems. Similarly, ProtoMotions~\cite{ProtoMotions} offers a comprehensive framework for physically plausible humanoid control and reframes character control as motion inpainting from partial goals.

In production XR stacks, humanoid rigs are often standardized and paired with engine-level retargeters and optional IK stages. This approach enables cross-asset compatibility but also introduces dependencies on specific SDKs or rig hierarchies. While SMPL-style models facilitate learning-based synthesis, real-time XR systems still prioritize lightweight solvers for responsiveness and low latency. Hand representation remains particularly error-prone, as mismatches in local frames or phalanx orientation often lead to perceptual artifacts unless anatomically consistent mappings are enforced. Thus, despite advances in body modeling and generative frameworks, robust runtime alignment remains a challenge.\\

% Ours
Unlike prior XR frameworks that couple motion to fixed pipelines and SDK-specific flows, MotionDAG decomposes motion processing into modular, event-driven nodes connected through a shared kinematic buffer. This design enables backend sensors to be swapped without refactoring, supports node-level ablations, and maintains consistent latency analysis. By normalizing heterogeneous inputs into a unified pose stream and applying device-agnostic retargeting, MotionDAG achieves stable, expressive embodiment across diverse rigs within a single runtime graph.